% !TEX program = xelatex
% ============================================================
%  Arabic Academic CV Template — Nibras Center
%  Compile with: xelatex main.tex (run twice)
%  Engine: XeLaTeX (required for polyglossia + bidi + fontspec)
% ============================================================

\documentclass[11pt,a4paper]{article}

% ---- Geometry ----
\usepackage[a4paper,margin=2cm,top=2cm,bottom=2cm]{geometry}

% ---- Core packages (order matters for bidi) ----
\usepackage{bidi}
\usepackage[no-math]{fontspec}
\usepackage[quiet]{polyglossia}

% ---- Fonts ----
\setmainfont[Scale=0.95]{Amiri-Regular.ttf}
\newfontfamily\arabicfont[Script=Arabic,Scale=0.95]{Amiri-Regular.ttf}
\newfontfamily\englishfont[Scale=0.85]{TeX Gyre Termes}

% ---- Languages ----
\setdefaultlanguage[calendar=gregorian,hijricorrection=1]{arabic}
\setotherlanguage[variant=british]{english}

% ---- Hyperref (load after polyglossia/bidi) ----
\usepackage[hidelinks]{hyperref}
\hypersetup{
  pdftitle={السيرة الذاتية الأكاديمية},
  pdfauthor={د. اسم الباحث}
}

% ---- Section / list / table packages ----
\usepackage{titlesec}
\usepackage{enumitem}
\usepackage{tabularx}
\usepackage{booktabs}
\usepackage{xcolor}
\usepackage{fontawesome5}
\usepackage{setspace}

% ---- Colors ----
\definecolor{accent}{HTML}{1F3A5F}
\definecolor{rule}{HTML}{555555}

% ---- Section formatting (colored, with rule underneath) ----
% \titleformat* formats the starred (unnumbered) version of a section.
\titleformat*{\section}{\Large\bfseries\color{accent}}
\titleformat*{\subsection}{\normalsize\bfseries\color{accent}}
\titlespacing*{\section}{0pt}{1.2ex plus 0.3ex}{0.6ex plus 0.2ex}
\titlespacing*{\subsection}{0pt}{0.8ex plus 0.2ex}{0.4ex plus 0.1ex}

% Rule under each \section* heading (redefine \section to handle the star)
\makeatletter
\let\oldsection\section
\renewcommand{\section}{\@ifstar{\sectionstar}{\oldsection}}
\newcommand{\sectionstar}[1]{\oldsection*{#1}\vspace{-0.6em}{\color{rule}\hrule height 0.4pt}\vspace{0.4em}}
\makeatother

% ---- List spacing (compact) ----
\setlist[itemize]{topsep=0.25em, itemsep=0.12em, parsep=0pt, leftmargin=*}
\setlist[enumerate]{topsep=0.25em, itemsep=0.12em, parsep=0pt, leftmargin=*}

% ---- Table column types ----
\newcolumntype{R}[1]{>{\raggedleft\arraybackslash}p{#1}}
\newcolumntype{C}[1]{>{\centering\arraybackslash}p{#1}}

% ---- Line spacing ----
\setstretch{1.1}

% ---- No page numbers on first page, plain elsewhere ----
\pagestyle{plain}

\begin{document}

% ============================================================
% HEADER
% ============================================================
\thispagestyle{empty}
\begin{center}
{\Huge\bfseries\color{accent} د. اسم الباحث\par}
\vspace{0.4em}
{\large أستاذ جامعي \,$\vert$\, باحث في التقنية العصبية\par}
\vspace{0.8em}
{\small
\faEnvelope~\href{mailto:name@example.com}{name@example.com} \quad
\faPhone~\LR{+90 555 000 0000} \quad
\faGlobe~\href{https://example.com}{example.com} \quad
\faMapMarker~\LR{Istanbul, Turkey} \quad
\faOrcid~\href{https://orcid.org/0000-0000-0000-0000}{\LR{ORCID}}
\par}
\end{center}

\vspace{0.5em}

% ============================================================
% 1. الملخص المهني
% ============================================================
\section*{الملخص المهني}
باحث جامعي متخصص في التقنية العصبية وواجهات الدماغ والحاسوب (\LR{BCI})،
مع خبرة تتجاوز خمسة عشر عاماً في التدريس والبحث العلمي. نشر أكثر من ثلاثين
ورقة بحثية في مجلات محكّمة، وأشرف على عدة مشاريع مموّلة وطُلّاب دراسات
عليا. يهتم بتطوير أنظمة مساعدة للأشخاص ذوي الإعاقة الحركية.

% ============================================================
% 2. التعليم
% ============================================================
\section*{التعليم}
\vspace{-0.3em}
\begin{tabularx}{\textwidth}{@{}l X c@{}}
\toprule
\textbf{الدرجة} & \textbf{المؤسسة — عنوان الأطروحة} & \textbf{السنة} \\
\midrule
دكتوراه & جامعة إسطنبول التقنية — \LR{Neural Decoding for BCI} & \LR{2012} \\
ماجستير & جامعة دمشق — تحليل الإشارات الحيوية & \LR{2008} \\
بكالوريوس & جامعة حلب — هندسة الطب الحيوي & \LR{2005} \\
\bottomrule
\end{tabularx}

% ============================================================
% 3. الخبرات الأكاديمية
% ============================================================
\section*{الخبرات الأكاديمية}
\textbf{أستاذ مشارك} — جامعة إسطنبول التقنية، قسم الهندسة الطبية الحيوية \hfill \LR{2018–الآن}\\
التدريس: معالجة الإشارات الحيوية، التعلم الآلي، الأخلاقيات البحثية. الإشراف على
رسائل الماجستير والدكتوراه.
\vspace{0.5em}

\textbf{أستاذ مساعد} — جامعة دمشق، قسم الهندسة الطبية الحيوية \hfill \LR{2012–2018}\\
تطوير مقررات معالجة الإشارات، والمشاركة في تأسيس مختبر واجهات الدماغ والحاسوب.

% ============================================================
% 4. المنشورات
% ============================================================
\section*{المنشورات}

\subsection*{كتب}
\begin{enumerate}
\item \textbf{مدخل إلى واجهات الدماغ والحاسوب.} دار النشر، \LR{2020}.
\item \textbf{معالجة الإشارات العصبية: الأسس والتطبيقات.} دار النشر، \LR{2017}.
\end{enumerate}

\subsection*{مقالات محكمة}
\begin{enumerate}
\item \LR{Author, A.} (2023). Deep learning for motor imagery \LR{BCI}. \emph{Journal of Neural Engineering}, \LR{20(3)}, 123--135.
\item \LR{Author, A.}, \LR{\& Coauthor, B.} (2021). Adaptive filters for \LR{EEG}. \emph{IEEE TNSRE}, \LR{29}, 1--10.
\end{enumerate}

\subsection*{مؤتمرات}
\begin{enumerate}
\item \LR{Author, A.} (2022). Real-time decoding of imagined speech. \emph{Proc. IEEE EMBC}, \LR{1--4}.
\item \LR{Author, A.}, \LR{\& Coauthor, B.} (2019). \LR{SSVEP}-based spelling system. \emph{Proc. ICONIP}, \LR{210--218}.
\end{enumerate}

% ============================================================
% 5. المشاريع البحثية
% ============================================================
\section*{المشاريع البحثية}
\textbf{نظام تواصل مساعد مبني على واجهات الدماغ والحاسوب} \hfill \LR{2021–2024}\\
\textit{الباحث الرئيسي} — تمويل: مجلس البحث العلمي (\LR{TUBITAK})، مبلغ \LR{300{,}000} دولار.

\vspace{0.4em}
\textbf{تطوير خوارزميات فك تشفير الكلام المتخيّل} \hfill \LR{2018–2020}\\
\textit{باحث مشارك} — تمويل: برنامج إطار البحث الأوروبي (\LR{Horizon 2020}).

% ============================================================
% 6. الجوائز والتكريمات
% ============================================================
\section*{الجوائز والتكريمات}
\begin{itemize}
\item جائزة البحث العلمي المتميز، جامعة إسطنبول التقنية، \LR{2022}.
\item زمالة ما بعد الدكتوراه، مؤسسة ألكسندر فون هومبولت، \LR{2015}.
\item أفضل ورقة بحثية شابة، مؤتمر \LR{ICONIP}، \LR{2019}.
\end{itemize}

% ============================================================
% 7. العضويات المهنية
% ============================================================
\section*{العضويات المهنية}
\begin{itemize}
\item الجمعية الدولية لهندسة الطب والأحياء (\LR{IEEE EMBS}) — عضو منذ \LR{2013}.
\item جمعية واجهات الدماغ والحاسوب العالمية (\LR{BCI Society}) — عضو منذ \LR{2016}.
\item هيئة التحرير، مجلة الهندسة العصبية العربية — \LR{2020–الآن}.
\end{itemize}

% ============================================================
% 8. اللغات
% ============================================================
\section*{اللغات}
\begin{itemize}
\item \textbf{العربية:} اللغة الأم.
\item \textbf{\LR{English}:} إجادة تامة (تحدثاً وكتابة).
\item \textbf{\LR{Turkish}:} إجادة جيدة.
\item \textbf{\LR{French}:} مستوى متوسط.
\end{itemize}

% ============================================================
% 9. المهارات التقنية
% ============================================================
\section*{المهارات التقنية}
\begin{itemize}
\item لغات البرمجة: \LR{Python}, \LR{R}, \LR{MATLAB}, \LR{C++}.
\item أدوات: \LR{LaTeX}, \LR{Git}, \LR{Docker}, \LR{Linux}.
\item تعلّم عميق: \LR{PyTorch}, \LR{TensorFlow}, \LR{Keras}.
\item معالجة الإشارات: \LR{MNE-Python}, \LR{EEGLAB}, \LR{FieldTrip}.
\end{itemize}

\end{document}
